\documentclass{article}
\usepackage{mathtools} 
\usepackage{fontspec}
\usepackage[UTF8]{ctex}
\usepackage{amsthm}
\usepackage{mdframed}
\usepackage{xcolor}
\usepackage{amssymb}
\usepackage{amsmath}


% 定义新的带灰色背景的说明环境 zremark
\newmdtheoremenv[
  backgroundcolor=gray!10,
  % 边框与背景一致，边框线会消失
  linecolor=gray!10
]{zremark}{说明}

% 通用矩阵命令: \flexmatrix{矩阵名}{元素符号}{行数}{列数}
\newcommand{\flexmatrix}[4]{
  \[
  #1 = \begin{pmatrix}
    #2_{11}     & #2_{12}     & \cdots & #2_{1#4}   \\
    #2_{21}     & #2_{22}     & \cdots & #2_{2#4}   \\
    \vdots      & \vdots      & \ddots & \vdots     \\
    #2_{#31}    & #2_{#32}    & \cdots & #2_{#3#4}
  \end{pmatrix}
  \]
}

% 简化版命令（默认矩阵名为A，元素符号为a）: \quickmatrix{行数}{列数}
\newcommand{\quickmatrix}[2]{\flexmatrix{A}{a}{#1}{#2}}

\begin{document}
\title{习题 17.1}
\author{张志聪}
\maketitle

\section*{7}

考虑
\begin{align*}
  \lim\limits_{(x, y) \to (0, 0)} f(x, y)
   & = \lim\limits_{(x, y) \to (0, 0)} (x^2 + y^2) sin \frac{1}{\sqrt{x^2 + y^2}} \\
   & \leq \lim\limits_{(x, y) \to (0, 0)} (x^2 + y^2) \frac{1}{\sqrt{x^2 + y^2}}  \\
   & = \lim\limits_{(x, y) \to (0, 0)} \sqrt{x^2 + y^2}                           \\
   & = 0 = f(0, 0)
\end{align*}
故$f(x)$在$(0, 0)$处连续。

\section*{15}

$f$在$U(P)$上的偏导数有界，故存在$M > 0$，对任意$(x, y) \in U(P)$，有
\begin{align*}
  |f_x(x, y)| < M \\
  |f_y(x, y)| < M
\end{align*}

对任意$(x_1, y_1) \in U(P)$，
由定理17.3（中值公式），我们有
\begin{align*}
   & \lim\limits_{(x, y) \to (x_1, y_1)} f(x, y) - f(x_1, y_1)                                                                   \\
   & = \lim\limits_{(x, y) \to (x_1, y_1)} f_x(x_1 + \theta_1(x - x_1), y)(x - x_1) + f_y(x_0, y_1 + \theta_2(y - y_1))(y - y_1) \\
   & < \lim\limits_{(x, y) \to (x_1, y_1)} M (x - x_1) + M(y - y_1)                                                              \\
   & = 0
\end{align*}
故
\begin{align*}
  \lim\limits_{(x, y) \to (x_1, y_1)} f(x, y) = f(x_1, y_1)
\end{align*}
命题得证。

\section*{16}

\begin{itemize}
  \item (1)

        $f(x, y)$的函数值，只与$y$有关，与$x$无关。

        对任意$(x_1, y), (x_2, y) \in int\, D$，由中值公式，我们有
        \begin{align*}
          f(x_1, y) - f(x_2, y)
           & = f_x(x_2 + \theta_1(x_1 - x_2), y) (x_1 - x_2) + f_y(x_2, y + \theta_2(y - y))(y - y) \\
           & = 0 + 0                                                                                \\
           & = 0
        \end{align*}
        由$(x_1, y), (x_2, y)$的任意性可知，在$int\, D$上，
        $f(x, y)$的函数值，只与$y$有关，与$x$无关。

        对任意$(x_0, y_0) \in \partial D$，构造一个元素都是$\int\, D$且$y = y_0$的点列${P_n}$收敛于$(x, y)$.
        由于$f$在$D$上连续，于是由定理16.5的推论可知$\{f(P_n)\}$收敛，设收敛于$A$，又因为$A$等于$\int\, D$中
        任意一个$y = y_0$的点，即与$x_0$无关。

        综上，命题得证。

  \item (2)

        $f(x, y)$在$int\, D$上常数函数，使用中值公式即可证明。

  \item (3)

        由证明过程可知，$f$在$D$上的连续性假设不能省略。比如
        \begin{equation*}
          f(x, y) =
          \begin{cases}
            0, & (a, b) \times (c, d)               \\
            1, & x = a, x = b \text{或} y = c, y = d
          \end{cases}
        \end{equation*}
\end{itemize}



\end{document}